% page 320 of the facsimile (image p-319 of the scan)
% block 162.6 x 246.3 mm ; left margin 8.0 mm
\pgc{-12.700mm}{0.000mm}{
\l{\cel{4}312}
\l{}
\l{}
\l{\sou{kratego}. (\sou{bot.}) I. Arbusto dornoza, ek la familio \textquotedbl{}rozacei\textquotedbl{},}
\l{\cel{3}kun mikra flori blanka (L.\cel{1}\sou{crataegus}).- II. La floro di}
\l{\cel{3}ta arbusto. - L.}
\l{}
\l{\sou{kratero}. I. (\sou{epoki antiqua}).Vazo larja, en qua onu mixis vino}
\l{\cel{3}ed aquo por plenigar la kalici e la kupi. - II. (\sou{metaf}.)}
\l{\cel{3}Aperturo formacita somite di volkano, per qua olca forjetas}
\l{\cel{3}la materii inflamita : flami, lavao, cindri. - III. Aperturo}
\l{\cel{3}facita ye la parto supera di fornego di vitrifisto. - DEFIRS}
\l{}
\l{\sou{kravacho}. Vergo flexebla qua finas per mecho, quan la kaval-}
\l{\cel{3}kanti uzas kom flogilo. - DF.}
\l{}
\l{\sou{kravato}. Peco de stofo quan onu +weras cirkum sua kolo, kom}
\l{\cel{3}ornivo, e qua (maxim-multa-kaze) pendas extremaje sur la}
\l{\cel{3}pektoro. - DEFI.}
\l{}
\l{\sou{krayono}. Peco de erco (grafito, e c.), plu o min mola, quan}
\l{\cel{3}onu uzas por trasar, desegnar, skribar, e c. - EF.}
\l{}
\l{\sou{krazo}. (\sou{gram.}) Fuzuro de la vokalo o diftongo finala di vor-}
\l{\cel{3}to, kun la vokalo o diftongo komencala di la vorto se-}
\l{\cel{3}quanta. - DEFIS.}
\l{}
\l{\sou{krear}. (\sou{trans., de, ek, per}). I.\cel{1}\sou{(religio}) (\sou{aludante deo}).}
\l{\cel{3}Facar (ulo) ek nulo, prenante lu de nihilo. - II. (\sou{aludant}e}
\l{\cel{3}\sou{homo)}\cel{1}. Kompozar (verko, kozo qua ne existis antee). -EFIS.}
\l{}
\l{\sou{krecento}. I. (1) Formo sesguroza di Luno dum lua kresko o des-}
\l{\cel{3}kresko ; (2) To quo havas ta formo. - II. Pano-bloko butro-}
\l{\cel{3}pasta, kun formo di Luno-krecento ; facata precipue en Paris,}
\l{\cel{3}Francia. - EFIS.}
\l{}
\l{\sou{kredar}. (\sou{trans.}). I. (1) Admisar(ulo) kom exakta sen verifiko.}
\l{\cel{3}(2) Egardar (ulu) kom dicanta lo exakta, sen verifikar lo.}
\l{\cel{3}- II. (ye ulu, ye ulo) : (1) Fidar a la realeso di ulo,sen}
\l{\cel{3}verifikar ; fidar a la efikiveso di ulo, ad olua povo, sen}
\l{\cel{3}verifikar ; (2) Fidar a la verdico di ulu, sen verifiko;}
\l{\cel{3}(3) Fidar a lua karaktero, sen verifiko. - III. (\sou{senco}\cel{1}ab-}
\l{\cel{3}\sou{soluta}). Fidar religiale. - EFIRS.}
\l{}
\l{\sou{kredenco.}\cel{2}Plado-moblo, konsolo o tablo sur qua onu pozas la}
\l{\cel{3}pladi, pladeti, glasi, e c. por la servi di tablo lor re-}
\l{\cel{3}pasto. - DEFIS.}
\l{}
\l{\sou{kreditar}. (\sou{trans., per ulo)}.\cel{2}Vendar, aceptante la ajorno di}
\l{\cel{3}la pago da la kompranto, pro ke onu judikas lu solventa.}
\l{\cel{3}- DEFIRS.}
\l{}
\l{\sou{kremo}. I. La parto maxim densa di lakto, qua acensas a la sur-}
\l{\cel{3}faco kande onu lasas lu repozar, ed ek qua onu facas butro}
\l{\cel{3}per quirlo. (\sou{metaf.)}\cel{1}La parto maxim bona de ulo. - II. Pe-}
\l{\cel{3}likulo qua eventas surface di lakto kande onu varmigas lu.-}
\l{\cel{3}III. Preparajo qua aspektas quale kremo. - DEFIRS.}
}
